\ProvidesFile{Logic.tex}[Логические операции]


\subsection{Логические операции}

\subsubsection{Вспомогательные классы}

Прежде всего я хочу уйти от ужасного синтаксиса STL вроде \verb"std::integral_constant" или \verb"std::conjunction".
Это кошмар.

\paragraph{Работа со значениями}

Для начала давайте определим базовые классы, с которыми мы будем работать.
В начале класс для произвольного типа, хранимого в compile-time.
\begin{cppcode}
template<class T, T x>
struct Value_t {
  using value_type = T;
  static constexpr value_type value = x;
};
\end{cppcode}
Здесь мы следуем соглашению, что класс-wrapper всегда хранит значение в статической переменной с именем \verb"value".
Псевдоним \verb"value_type" вводится для совместимости с STL.

Теперь можно задавать константы во время компиляции так
\begin{cppcode}
using Value1 = Value_t<int, 1>;
using Value2 = Value_t<bool, false>;
using Value3 = Value_t<char, 'c'>;
\end{cppcode}
Для удобства можно ввести вспомогательные типы вроде
\begin{cppcode}
template<bool v>
struct Bool_t : Value_t<bool, v> {};

template<char v>
struct Char_t : Value_t<char, v> {};

template<int v>
struct Int_t : Value_t<int, v> {};
...
\end{cppcode}
и так далее.
Можно создать соответствующие типы для знаковых и беззнаковых встроенных типов.
Так же можно создать соответствующие типы для \verb"float" и \verb"double".
И теперь можно писать
\begin{cppcode}
using Value1 = Int_t<1>;
using Value2 = Bool_t<false>;
using Value3 = Char_t<'c'>;
using Value4 = Float_t<1.2f>;
using Value5 = Double_t<2.>;
\end{cppcode}


\paragraph{Работа с типами}

В дальнейшем нам понадобится возможность упаковать внутрь класса другой класс.
Для этого тоже заведем класс-wrapper.
\begin{cppcode}
template<class T>
struct Type_t {
  using type = T;
};
\end{cppcode}
Еще понадобится возможность установить тип переменной.
Для этого можно использовать встроенный в язык \verb"decltype".
Однако, у него есть одна небольшая особенность, которую я хочу обсудить.
Например, если мы напишем код вида
\begin{cppcode}
struct A {
  using value_type = int;
  static constexpr value_type value = 1;
};
\end{cppcode}
Тогда
\begin{cppcode}
std::is_same_v<decltype(A::value), int>; // false
std::is_same_v<decltype(A::value), const int>; // true
\end{cppcode}
То есть в тип переменной у нас добавляет ненужное \verb"const".
В run-time это пожалуй разумно, потому что мы хотим отличать константные переменные от не константных.
Но в compile-time у нас все переменные константные и такое поведение может привести к неприятным проблемам с кодом.
Чтобы избавиться от появления \verb"const" в типе \verb"value" введем вспомогательный шаблон для определения типа так:
\begin{cppcode}
namespace detail {
template<auto E>
struct TypeOfImpl : Type_t<decltype(E)> {};
} // namespace detail

template<auto E>
using TypeOf = detail::TypeOfImpl<E>::type;
\end{cppcode}
\begin{cppcode}
std::is_same_v<TypeOf<A::value>, int>; // true
std::is_same_v<TypeOf<A::value>, const int>; // false
\end{cppcode}
Причина такого поведения следующая.
При подстановке \verb"A::value" в шаблон \verb"TypeOf" мы присваиваем к переменной определенной с типом \verb"auto".
В таком случае \verb"auto" забывает про \verb"const".
Это связано с тем, как работает в плюсах присваивание по значению.
Например
\begin{cppcode}
int x = 1;
const int y = 2;
int z;
z = x; // OK
z = y; // OK
\end{cppcode}
Мы можем положить в неконстантную переменную как неконстантное значение, так и константное.
Оно же просто скопируется и копия уже будет неконстантной.
Кроме того,
\begin{cppcode}
const int x = 1;
auto z = x;  // decltype(z) = int
// or
const auto z = x; // decltype(z) = const int
\end{cppcode}
Обратите внимание, что \verb"auto" не становится константной, даже если переменная справа от равенства константная.
Если вы хотите сделать \verb"z" константной, то надо явно добавить квалификатор \verb"const".
Можно запомнить так: \verb"auto" никогда не \verb"const" и никогда не \verb"reference".

\subsubsection{Базовые связки}

Теперь нужно научиться манипулировать логическими переменными.
Например нам может понадобиться связка \verb"IfElse", чтобы выбрать тип в зависимости от условия.
Данная деятельность не является необходимой, но я хочу уйти от ужасного синтаксиса STL и дополнительно добавлю недостающие связки.

\paragraph{Значения и классы-wrapper-ы}

Прежде чем говорить про код, давайте обратим внимание, что у нас есть переменная и переменная запакованная в класс-wrapper:
\begin{cppcode}
constexpr int v = 1;

struct A {
  static constexpr int value = 2;
};
\end{cppcode}
Таким образом есть значения типа \verb"T", а есть эти значения помещенные в какой-то класс-wrapper \verb"W(T)".
Все логические операторы в STL работают над \verb"W(bool)" и возвращают значение из \verb"W(bool)".
В общем случае у нас есть $4$ разных варианта:
\begin{enumerate}
\item \verb"T"$\to$\verb"T".
Принимаем значения и возвращаем значения.

\item \verb"T"$\to$\verb"W(T)".
Принимаем значения, результат возвращаем завернутым в класс.
Я для такого случая буду использовать суффикс \verb"_t".

\item \verb"W(T)"$\to$\verb"W(T)".
Принимаем аргументы завернутыми в класс и возвращаем результат завернутым в класс.
Для такого случая я буду использовать суффикс \verb"_w".

\item \verb"W(T)"$\to$\verb"T".
Принимаем аргументы завернутыми в класс и возвращаем значение.
Для такого случая я буду использовать суффикс \verb"_u".
\end{enumerate}

\paragraph{Имплементация связок}

Вот пример того, как это выглядит для оператора \verb"And":
\begin{cppcode}
namespace detail {
template<class /*void*/, class W1, class... Ws>
struct AndImpl : W1 {};

template<class W1, class W2, class... Ws>
struct AndImpl<std::enable_if_t<bool(W1::value)>, W1, W2, Ws...>
    : AndImpl<void, W2, Ws...> {};
} // namespace detail

template<class... Ws>
struct And_w : detail::AndImpl<void, Ws...> {};

template<>
struct And_w<> : Bool_t<true> {};

template<class... Ws>
constexpr bool And_u = And_w<Ws...>::value;

template<bool... Cs>
struct And_t : And_w<Bool_t<Cs>...> {};

template<bool... Cs>
constexpr bool And = And_t<Cs...>::value;
\end{cppcode}
Зачем нужны все $4$ случая я обсужу чуть позже, а сейчас хочу обратить внимание на то, как работает эта имплементация.
Для имплементации я завожу вспомогательный аргумент-флаг.
Вторая версия шаблона является более специализированной версией первого.
Потому компилятор будет стараться использовать его всегда, когда можно.
Вторая специализация используется только если внутри первого аргумента \verb"W1" истина.
Если же внутри \verb"W1" лежит ложь, то шаблон хранит значение \verb"value" из \verb"W1", то есть ложь.
Обратите внимание, что если \verb"W1" хранит ложь, то остальные аргументы не вычисляются.
Это очень важное явление, которое нам пригодится в будущем.

Аналогично можно определить и все остальные операции.
Вот так можно сделать операцию логического <<или>>:
\begin{cppcode}
namespace detail {

template<class /*void*/, class W1, class... Ws>
struct OrImpl : W1 {};

template<class W1, class W2, class... Ws>
struct OrImpl<std::enable_if_t<!bool(W1::value)>, W1, W2, Ws...>
    : OrImpl<void, W2, Ws...> {};
} // namespace detail

template<class... Ws>
struct Or_w : detail::OrImpl<void, Ws...> {};

template<>
struct Or_w<> : Bool_t<false> {};

template<class... Ws>
constexpr bool Or_u = Or_w<Ws...>::value;

template<bool... Cs>
struct Or_t : Or_w<Bool_t<Cs>...> {};

template<bool... Cs>
constexpr bool Or = Or_t<Cs...>::value;
\end{cppcode}
Отрицание можно сделать так:
\begin{cppcode}
namespace detail {

template<class W>
struct NotImpl : Bool_t<!bool(W::value)> {};
} // namespace detail

template<class W>
struct Not_w : detail::NotImpl<W> {};

template<class W>
constexpr bool Not_u = Not_w<W>::value;

template<bool C>
struct Not_t : Not_w<Bool_t<C>> {};

template<bool C>
constexpr bool Not = Not_t<C>::value;
\end{cppcode}
Еще добавим импликацию
\begin{cppcode}
namespace detail {
template<class W1, class W2>
struct ImplyImpl : Or_w<Not_w<W1>, W2> {};
}; // namespace detail

// A => B
template<class W1, class W2>
struct Imply_w : detail::ImplyImpl<W1, W2> {};

template<class W1, class W2>
constexpr bool Imply_u = Imply_w<W1, W2>::value;

template<bool A, bool B>
struct Imply_t : Imply_w<Bool_t<A>, Bool_t<B>> {};

template<bool A, bool B>
constexpr bool Imply = Imply_t<A, B>::value;
\end{cppcode}
Теперь обсудим как пишется \verb"IfElse".
\begin{cppcode}
namespace detail {

template<bool condition, class W1, class W2>
struct IfElseImpl : W2 {};

template<class W1, class W2>
struct IfElseImpl<true, W1, W2> : W1 {};
} // namespace detail

template<bool condition, class W1, class W2>
struct IfElse_w : detail::IfElseImpl<condition, W1, W2> {};

template<bool condition, class W1, class W2>
constexpr auto IfElse_v = IfElse_w<condition, W1, W2>::value;

template<bool condition, class W1, class W2>
using IfElse_u = typename IfElse_w<condition, W1, W2>::type;

template<bool condition, class T1, class T2>
struct IfElse_t : IfElse_w<condition, Type_t<T1>, Type_t<T2>> {};

template<bool condition, class T1, class T2>
using IfElse = typename IfElse_t<condition, T1, T2>::type;
\end{cppcode}

\paragraph{Еще связки}

Я так же определяю еще связки для сравнения разного вида, чтобы можно было сравнивать сразу несколько значений на равенство друг другу.
\begin{cppcode}
namespace detail {
template<class W1, class W2>
struct areSame2Impl : Bool_t<false> {};

template<class W>
struct areSame2Impl<W, W> : Bool_t<true> {};
} // namespace detail

template<class T1, class T2>
struct areSame2_t : detail::areSame2Impl<T1, T2> {};

template<class T1, class T2>
constexpr bool areSame2 = areSame2_t<T1, T2>::value;

template<class W1, class W2>
struct areSame2_w : areSame2_t<typename W1::type, typename W2::type> {};

template<class W1, class W2>
constexpr bool areSame2_u = areSame2_w<W1, W2>::value;

namespace detail {

template<class... Ws>
struct areSameImpl : Bool_t<false> {};

template<class W1, class W2, class... Ws>
struct areSameImpl<W1, W2, Ws...>
    : And_w<areSame2_w<W1, W2>, areSameImpl<W2, Ws...>> {};

template<class W1>
struct areSameImpl<W1> : Bool_t<true> {};

template<>
struct areSameImpl<> : Bool_t<true> {};
} // namespace detail

template<class... Ws>
using areSame_w = detail::areSameImpl<Ws...>;

template<class... Ws>
constexpr bool areSame_u = areSame_w<Ws...>::value;

template<class... Ts>
struct areSame_t : detail::areSameImpl<Type_t<Ts>...> {};

template<class... Ts>
constexpr bool areSame = areSame_t<Ts...>::value;
\end{cppcode}
Так же определяю связки:
\begin{multicols}{2}
\begin{enumerate}
\item \verb"isLess"
\item \verb"isNotLess"
\item \verb"isGreater"
\item \verb"isNotGreater"
\item \verb"isEqual"
\item[]
\end{enumerate}
\end{multicols}
Приведу имплементацию двух связок.
Первая -- строгое сравнение \verb"isLess":
\begin{cppcode}
namespace detail {
template<class W1, class W2>
struct isLessImpl : Bool_t<(W1::value < W2::value)> {};
} // namespace detail

template<class W1, class W2>
struct isLess_w : detail::isLessImpl<W1, W2> {};

template<class W1, class W2>
constexpr bool isLess_u = isLess_w<W1, W2>::value;

template<auto A, auto B>
struct isLess_t
    : isLess_w<Value_t<Values::TypeOf<A>, A>, Value_t<Values::TypeOf<B>, B>> {};

template<auto A, auto B>
constexpr bool isLess = isLess_t<A, B>::value;
\end{cppcode}
Вторая связка -- сравнение на равенство:
\begin{cppcode}
template<class W1, class W2>
struct isEqualImpl : Bool_t<(W1::value == W2::value)> {};
} // namespace detail

template<class W1, class W2>
struct isEqual_w : detail::isEqualImpl<W1, W2> {};

template<class W1, class W2>
constexpr bool isEqual_u = isEqual_w<W1, W2>::value;

template<auto A, auto B>
struct isEqual_t
    : isEqual_w<Value_t<Values::TypeOf<A>, A>, Value_t<Values::TypeOf<B>, B>> {
};

template<auto A, auto B>
constexpr bool isEqual = isEqual_t<A, B>::value;
\end{cppcode}
Здесь полезными являются версии, которые принимают не сами значения, а их классы-wrapper-ы, потому что они позволяют сделать сравнение <<не вынимая значение из класса>>.
Однако, я всегда определяю все $4$ версии шаблонов.

\subsubsection{Условные вычисления}
\label{section::ConditionalComputation}

В этом разделе я хочу рассказать еще одну причину, зачем многие операции заворачиваются во вспомогательный класс.
Давайте начнем с примера.
Пусть у нас есть следующий шаблон для работы со списком
\begin{cppcode}
template<auto... Es>
struct List {};

using List1 = List<1, 2., 'c'>;
using List2 = List<>;
\end{cppcode}
Теперь имплементируем функцию проверки, что у нас список
\begin{cppcode}
namespace detail {

template<class TList>
struct isListImpl : Bool_t<false> {};

template<template<auto...> class TList, auto... Es>
struct isListImpl<TList<Es...>> : Bool_t<true> {};
} // namespace detail

template<class TList>
struct isList_t : detail::isListImpl<TList> {};

template<class TList>
constexpr bool isList = isList_t<TList>::value;

template<class WList>
struct isList_w : isList_t<typename WList::type> {};

template<class WList>
constexpr bool isList_u = isList_w<WList>::value;
\end{cppcode}
И еще добавим функцию вычисления размера списка так:
\begin{cppcode}
namespace detail {
template<class TList>
struct SizeImpl {
  static_assert(isList<TList>, "The argument of Size MUST be an element list!");
};

template<template<auto...> class TList, auto... Es>
struct SizeImpl<TList<Es...>> : Int_t<sizeof...(Es)> {};
} // namespace detail

template<class TList>
struct Size_t : detail::SizeImpl<TList> {};

template<class TList>
constexpr int Size = Size_t<TList>::value;

template<class WList>
struct Size_w : Size_t<typename WList::type> {};

template<class WList>
constexpr int Size_u = Size_w<WList>::value;
\end{cppcode}
Теперь давайте обсудим как написать проверку того, что список не пуст.
Наивный подход будет такой:
\begin{cppcode}
template<class TList>
constexpr bool isNonEmptyList = isList<TList> && (Size<TList> > 0);
\end{cppcode}
Однако, это не сработает.
В случае если \verb"TList" не является списком, то выражение \verb"Size<TList>" упадет на \verb"static_assert" из его имплементации.
То есть во время компиляции оператор \verb"&&" НЕ вычисляется лениво, должны быть вычислены все его аргументы.
Тем не менее, выше мы видели что на уровне классов-wrapper-ов логические связки <<и>> и <<или>> работают лениво.
Давайте введем еще функцию проверки на пустоту:
\begin{cppcode}
namespace detail {

template<class TList, bool isCorrect = isList<TList>>
struct isEmptyImpl {
  static_assert(isList<TList>,
                "The first argument of isEmpty MUST be an element list!");
};

template<class TList>
struct isEmptyImpl<TList, true> : Bool_t<(Size<TList> == 0)> {};
} // namespace detail

template<class TList>
struct isEmpty_t : detail::isEmptyImpl<TList> {};

template<class TList>
constexpr bool isEmpty = isEmpty_t<TList>::value;

template<class WList>
struct isEmpty_w : isEmpty_t<typename WList::type> {};

template<class WList>
constexpr bool isEmpty_u = isEmpty_w<WList>::value;
\end{cppcode}
Вот теперь мы можем сделать вот такую проверку на непустой список:
\begin{cppcode}
template<class TList>
struct isNonEmptyList_t : And_w<isList_t<TList>, Not_w<isEmpty_t<TList>>> {};

template<class TList>
constexpr bool isNonEmptyList = isNonEmptyList_t<TList>::value;

template<class WList>
struct isNonEmptyList_w : isNonEmptyList_t<typename WList::type> {};

template<class WList>
constexpr bool isNonEmptyList_u = isNonEmptyList_w<WList>::value;
\end{cppcode}
Тут нас будет больше всего интересовать первая версия с суффиксом \verb"_t".
\begin{cppcode}
template<class TList>
struct isNonEmptyList_t : And_w<isList_t<TList>, Not_w<isEmpty_t<TList>>> {};
\end{cppcode}
Здесь компилятор вычисляет два класса:
\begin{enumerate}
\item \verb"isList_t<TList>".
Этот шаблон вычисляется всегда по определению.

\item \verb"Not_w<isEmpty_t<TList>>".
Этот шаблон тоже вычисляется всегда.
Действительно \verb"isEmpty_t" существует для любого значения \verb"TList".
Он НЕ содержит \verb"value" внутри себя, а потому обращение внутрь этого класса будет вызывать ошибку компиляции и падение по \verb"static_assert".
Но если мы не будем залазить внутрь, то ошибки компиляции не будет.
\end{enumerate}
и потом подставляет их в качестве аргументов в шаблон \verb"And_w", который принимает классы-wrapper-ы и возвращает тоже класс-wrapper.
А так как мы унаследовались от него, то теперь и \verb"isNonEmptyList_t" будет содержать статическую константу \verb"value" с нужным значением.

\subsubsection{Проверки}
\label{section::LogicChecks}

Следующий этап -- умение делать проверки со вспомогательными типами.
В частности нужно уметь проверять есть ли в классе статическая переменная \verb"value" или тип \verb"type".
Так же полезно уметь проверять тип \verb"value".
Начнем с простого случая: проверяем, что элемент имеет заданный тип:
\begin{cppcode}
namespace detail {
template<auto v, class T>
struct isOfTypeImpl : std::is_same<TypeOf<v>, std::remove_cv_t<T>> {};
} // namespace detail

template<auto v, class T>
constexpr bool isOfType = detal::isOfTypeImpl<v, T>::value;
\end{cppcode}
Обратите внимание, что мы убираем \verb"const" и \verb"volatile" с типа \verb"T".

Теперь обсудим, как устроена проверка наличия полей \verb"value" и \verb"type" в классе.
Идея в том, чтобы обратиться к этим полям.
Если произошла ошибка, то поля нет, если ошибка не произошла, то поле есть.
Однако, мы не хотим, чтобы код падал с ошибкой.
Для этого вспоминаем SFINAE.
По-простому шаблонам разрешено иметь ошибки при специализации шаблона.
В этом случае это НЕ вызывает ошибку компиляции, а компилятор просто отбрасывает данную специализацию, как неподходящую.
Для этого вводится дополнительный параметр, в котором будет вычисляться нужное выражение.
Но прежде всего нам понадобится вспомогательный класс, который забывает, что в него подставили:
\begin{cppcode}
namespace detail {
template<class... T>
struct ForgetTImpl : Type_t<void>, Int_t<0> {};

template<auto... v>
struct ForgetVImpl : Type_t<void>, Int_t<0> {};
} // namespace detail

template<class... T>
using ForgetT = typename detail::ForgetTImpl<T...>::type;

template<auto... v>
using ForgetV = typename ForgetVImpl<v...>::type;
\end{cppcode}
Работает это так
\begin{cppcode}
using Type1 = ForgetT<int, double, char>; // Type1 = void
using Type2 = ForgetV<1, 2., 'c'>; // Type2 = void
\end{cppcode}
То есть этот шаблон всегда возаращает \verb"void".
Его польза в том, что в него можно подставить любые типы или любые выражения и вычислить их в качестве аргументов.
И вот этот шаблон и будет тем местом, где будет выполняться обращение к предполагаемому ресурсу.
Вот Как можно проверить наличие поля \verb"type" в класс:
\begin{cppcode}
namespace detail {

template<class T, class Flag = void>
struct hasTypeImpl : Bool_t<false> {};

template<class T>
struct hasTypeImpl<T, ForgetT<typename T::type>> : Bool_t<true> {};
} // namespace detail

template<class T>
constexpr bool hasType = detail::hasTypeImpl<T>::value;
\end{cppcode}
Давайте обсудим, как это работает.
Пусть у нас есть следующий код
\begin{cppcode}
struct A {};
struct B {
  using type = int;
};

std::cout << hasType<A> << '\n'; // false
std::cout << hasType<B> << '\n';  // true
\end{cppcode}
Обсудим, как работает код в случае класса \verb"A".
\begin{enumerate}
\item Для вычисления константы \verb"hasType" компилятор вычисляет правую часть, а именно
\begin{center}
\verb"hasTypeImpl<A, void>".
\end{center}

\item Теперь мы определили параметры для шаблона и надо понять, какую из двух версий использовать.
\begin{enumerate}
\item Под первую дефолтную версию подходят любые наборы параметров и получаем
\begin{center}
\verb"hasTypeImpl<A, void>".
\end{center}

\item Для второй специализации, компилятор пытается вычислить выражение
\begin{center}
\verb"hasTypeImpl<A, ForgetT<typename A::type>>".
\end{center}
Так как в \verb"A" нет поля \verb"type", то при попытке вычислить аргумент для \verb"ForgetT" будет ошибка.
А это значит, данная специализация не подходит.
И потому выбирается первая специализация.
\end{enumerate}

\item По определению в первой специализации значение \verb"value" равно \verb"false".
\end{enumerate}
Теперь посмотрим, как это работает в случае класса \verb"B".
\begin{enumerate}
\item Для вычисления константы \verb"hasType" компилятор вычисляет правую часть, а именно
\begin{center}
\verb"hasTypeImpl<B, void>".
\end{center}

\item Теперь мы определили параметры для шаблона и надо понять, какую из двух версий использовать.
\begin{enumerate}
\item Под первую дефолтную версию подходят любые наборы параметров и в частности \verb"<B, void>".

\item Для второй специализации, компилятор пытается вычислить выражение
\begin{center}
\verb"hasTypeImpl<B, ForgetT<typename B::type>>".
\end{center}

В итоге опять получаем \verb"hasTypeImpl<B, void>".
И эта специализация тоже подходит.
\end{enumerate}

\item Теперь компилятор должен выбрать какую версию шаблона взять.
Тут действует правило, что если одна версия шаблона является специализацией другой, то специализация всегда в приоритете.
В этом случае вторая версия является более специализированной и потому выбирается она.

\item По определению во второй специализации значение \verb"value" равно \verb"true".
\end{enumerate}
Этот же подход можно использовать для остальных проверок.
Вот проверка на наличие статической переменной \verb"value":
\begin{cppcode}
namespace detail {

template<class T, class Flag = void>
struct hasValueImpl : Bool_t<false> {};

template<class T>
struct hasValueImpl<T, ForgetV<T::value>> : Bool_t<true> {};
} // namespace detail

template<class T>
constexpr bool hasValue = detail::hasValueImpl<T>::value;
\end{cppcode}
Если же мы хотим проверить, что \verb"value" имеет специальный тип, то надо сделать так:
\begin{cppcode}
namespace detail {
template<class T, class R, class Flag = void>
struct hasValueOfTypeImpl : Bool_t<false> {};

template<class T, class R>
struct hasValueOfTypeImpl<T, R, ForgetV<T::value>> : Bool_t<isOfType<T::value, R>> {};
} // namespace detail

template<class T, class R>
constexpr bool hasValueOfType = detail::hasValueOfTypeImpl<T, R>::value;
\end{cppcode}
В таком случае, если аргументы шаблона дали ошибку, то есть нет поля \verb"T::value", то выражение проверяющее тип даже не будет вычисляться.
А если такое поле есть, то мы можем вызвать \verb"isOfType<T::value, R>".
И теперь можно написать проверку, что у нас есть поля \verb"value" и \verb"value_type" и что \verb"value" имеет тип \verb"value_type".
Делаем это так же через дополнительное поле:
\begin{cppcode}
namespace detail {
template<class T, class Flag = void>
struct isValueImpl : Bool_t<false> {};

template<class T>
struct isValueImpl<T, ForgetT<decltype(T::value), typename T::value_type>>
    : isOfType_t<T::value, typename T::value_type> {};
} // namespace detail

template<class T>
constexpr bool isValue = detail::isValueImpl<T>::value;
\end{cppcode}
А теперь добавим проверку, что \verb"value_type" совпадает с заданным типом \verb"R".
Делаем так:
\begin{cppcode}
namespace detail {
template<class T, class R, class Flag = void>
struct isValueOfTypeImpl : Bool_t<false> {};

template<class T, class R>
struct isValueOfTypeImpl<T, R,
                         ForgetT<decltype(T::value), typename T::value_type>>
    : Bool_t<isOfType<T::value, typename T::value_type> &&
             std::is_same_v<std::remove_cv_t<typename T::value_type>,
                            std::remove_cv_t<R>>> {};
} // namespace detail

template<class T>
constexpr bool isBool = detail::isValueOfTypeImpl<T, bool>::value;

template<class T>
constexpr bool isInt = detail::isValueOfTypeImpl<T, int>::value;
\end{cppcode}

\paragraph{Имя переменной}
% TO DO
Что если мы хотим проверить наличие не имени \verb"value" или \verb"type", а какого-то произвольного?
Вот тут у меня только плохие новости.
Механизмов изменять имя переменной в языке нет.
То есть под каждую переменную придется писать очередную проверку снова и снова.
И вот это то самое единственное место, где могут пригодиться макросы.
Надо превратить вот этот код в макрос:
\begin{cppcode}
namespace detail {
template<class T, class R, class Flag = void>
struct hasValueOfTypeImpl : Bool_t<false> {};

template<class T, class R>
struct hasValueOfTypeImpl<T, R, ForgetV<T::value>> : Bool_t<isOfType<T::value, R>> {};
} // namespace detail

template<class T, class R>
struct hasValueOfType_t : detail::hasValueOfTypeImpl<T, R> {};

template<class T, class R>
constexpr bool hasValueOfType = hasValueOfType_t<T, R>::value;

template<class WT, class WR>
struct hasValueOfType_w
    : hasValueOfType_t<typename WT::type, typename WR::type> {};

template<class WT, class WR>
constexpr bool hasValueOfType_u = hasValueOfType_w<WT, WR>::value;
\end{cppcode}
В начале сделаем вспомогательный макрос для генерации имплементации в пространстве имен \verb"detail":%
\footnote{Операция \texttt{\#\#} для препроцессора является конкатенацией строк.}
\begin{cppcode}
#define HAS_MEMBER_OF_TYPE_IMPL(member)                                        \
  namespace detail {                                                           \
  template<class T, class R, class Flag = void>                                \
  struct has_##member##_OfTypeImpl : Bool_t<false> {};                         \
  template<class T, class R>                                                   \
  struct has_##member##_OfTypeImpl<T, R, ForgetV<T::member>>                   \
      : isOfType_t<T::member, R> {};                                           \
  }
\end{cppcode}
А теперь сам макрос
\begin{cppcode}
#define HAS_MEMBER_OF_TYPE(member)                                             \
  HAS_MEMBER_OF_TYPE_IMPL(member)                                              \
  template<class T, class R>                                                   \
  struct has_##member##_OfType_t : detail::has_##member##_OfTypeImpl<T, R> {}; \
  template<class T, class R>                                                   \
  constexpr bool has_##member##_OfType = has_##member##_OfType_t<T, R>::value; \
  template<class WT, class WR>                                                 \
  struct has_##member##_OfType_w                                               \
      : has_##member##_OfType_t<typename WT::type, typename WR::type> {};      \
  template<class WT, class WR>                                                 \
  constexpr bool has_##member##_OfType_u =                                     \
      has_##member##_OfType_w<WT, WR>::value;
\end{cppcode}
Теперь надо в начале сгенерировать шаблон для нужного имени и типа:
\begin{cppcode}
HAS_MEMBER_OF_TYPE(x)                                            
\end{cppcode}
Этот шаблон развернется в
\begin{cppcode}
namespace detail {
template<class T, class R, class Flag = void>
struct has_x_OfTypeImpl : Bool_t<false> {};

template<class T, class R>
struct has_x_OfTypeImpl<T, R, ForgetV<T::x>> : isOfType_t<T::x, R> {};
}

template<class T, class R>
struct has_x_OfType_t : detail::has_x_OfTypeImpl<T, R> {};

template<class T, class R>
constexpr bool has_x_OfType = has_x_OfType_t<T, R>::value;

template<class WT, class WR>
struct has_x_OfType_w : has_x_OfType_t<typename WT::type, typename WR::type> {};

template<class WT, class WR>
constexpr bool has_x_OfType_u = has_x_OfType_w<WT, WR>::value;
\end{cppcode}
И теперь его можно использовать так:
\begin{cppcode}
class A {
public:
  static constexpr int x = 0;
};

class B {
  static constexpr int x = 0;
};

struct C {
  static constexpr bool x = false;
};

int main() {
  std::cout << "has int x " << has_x_OfType<A, int> << '\n';  // returns true
  std::cout << "has int x " << has_x_OfType<B, int> << '\n';  // returns false
  std::cout << "has int x " << has_x_OfType<C, int> << '\n';  // returns false
  return 0;
}
\end{cppcode}
Обратите внимание, что тут идет проверка на доступную статическую переменную типа \verb"int" с именем \verb"x" внутри класса.
Правильная проверка на наличие тех или иных полей в классе -- это отдельный вопрос.

\subsubsection{Операции}

Непосредственно над типами операции можно определять с помощью шаблонов.
Однако, это автоматически означает, что вы лишаетесь возможность использовать привычные операторы для работы со значениями, например, сложить два целых числа или склеить две строки.
Однако, можно вместо этого перейти к переменным и перегрузить все операции.
Вот примеры
\begin{cppcode}
template<class T, T x>
constexpr Value_t<T, x> val_v = {};

template<bool v>
constexpr Bool_t<v> bool_v = {};

template<int v>
constexpr Int_t<v> int_v = {};
\end{cppcode}
Теперь можно определить операции, вот например сложение.
\begin{cppcode}
namespace detail {
template<class T>
concept ClassValue = isValue<T>;
} // namespace detail

template<detail::ClassValue A, detail::ClassValue B>
constexpr auto operator+(A, B) {
  using T = decltype(A::value + B::value);
  return val_v<T, A::value + B::value>;
}
\end{cppcode}
Тогда можно использовать следующий синтаксис
\begin{cppcode}
auto value1 = int_v<1> + int_v<2>;
// value1 = val_v<int, 3>;
\end{cppcode}
Для удобства можно добавить конвертацию переменных к хранимому типу.
Я не люблю делать неявный каст, потому сделаю оператор скобки
\begin{cppcode}
template<class T, T x>
struct Value_t {
  using value_type = T;
  static constexpr value_type value = x;

  constexpr value_type operator()() const noexcept {
    return value;
  }
};
\end{cppcode}
Тогда можно пользоваться так
\begin{cppcode}
auto value1 = int_v<1> + int_v<2>;
// value1 = val_v<int, 3>;
std::cout << "value1 = " << value1() << '\n';
\end{cppcode}
Аналогично можно перегрузить все существующие операторы (бинарные и унарные), включая операторы сравнения.
